\chapter*{General Introduction}
\phantomsection
\addcontentsline{toc}{chapter}{General Introduction}
\label{chap:cascade-general-introduction}

The internship comprised several related contributions, but not every assignment fed
directly into the benchmark campaigns. The first documented task was an independent board
equivalence search: the host team provided a list of competitor boards and asked us to find
their nearest STM32 counterparts. This product-comparison deliverable did not determine the
boards or scenarios used later in the project.

The experimental work began separately with a pilot Universal Serial Bus (USB) device
campaign on ST, Renesas, and NXP platforms. The pilot combined footprint and initialisation
measurements with an optimisation pass on the STM32 reference projects. It also exposed
repeated manual work in map-file reconciliation, project preparation, and evidence
tracking. We developed MCU Workflow Studio and the MCU Benchmark Copilot Agent in response
to those limitations, then applied the broader methodology in the later GD32 and Texas
Instruments campaigns.

The objective is not to produce a universal vendor ranking. Processor architecture,
memory organisation, clock configuration, compiler output, middleware boundaries, and
scenario design all affect the result. We consequently compare ST with one competing
vendor at a time and restrict each conclusion to the recorded board pair, configuration,
and scenario. Where internal workbook cells disagree with reviewed result images, the
report uses the displayed presentation values and records that source decision explicitly.

Chapter~\ref{chap:general-description} presents the industrial context and requirements.
Chapter~\ref{chap:cascade-board-equivalence} documents the standalone board-equivalence
assignment. Chapter~\ref{chap:cascade-usb-pilot} then begins the independent experimental
work. The two
supporting-tool chapters follow because their requirements emerged from that pilot. The
GD32 and TI chapters then show how the work expanded beyond the initial USB studies, and
the conclusion returns to the limits that govern all of these comparisons.